% This file should be rename with the speaker's last name. (e.g : smith)

% Write the author names in the order you wish them to appear
% and underline the one who will give the talk.

% Only the speaker's email address is required.

\begin{abstract_online}{Realizing the concept of “Multiple Representations” by using CAS (Part I, Part II)}{%
        \underline{Helmut Heugl}$^{1}$}{[hheugl@aon.at]
        }{%
        $^1$ Head of ACDCA (Austrian Center for Didactics of Computer Algebra)}
    
Mathematical concepts are presented in multiple modes of representation (or “prototypes”) such as text, graphs and diagrams, tables, algebraic expressions and computer simulations.  A prime goal of teaching is to help learners develop an understanding of the mathematical concepts by considering and using these different representational modes and levels. Several prototypes of the concept provide complementary information [1]. 
Therefore it is not enough to become acquainted with and to understand the information of a certain representation mode. A central cognitive activity on the way to mathematical concepts is to build links between representation modes of a concept. In traditional mathematics education prototypes mostly are available in a serial way.  The main importance of technology tools is that the learner can use several prototypes parallely.
By using examples of Algebra and Analysis I will show the role of CAS when building links between several representation modes of a concept or when solving problems [2].




\textbf{References} \newline{}
% Model for references to books.
[1] \textsc{Dörfler, W.}, Der Computer als kognitives Werkzeug und kognitives Medium.  In \emph{Computer - Mensch – Mathematik. Verlag Hölder-Pichler-Tempsky}, Wien, 1991, pp51. ISBN3-209-01452. \newline{}
%  Model for references to articles
[2] \textsc{Heugl, H.}, Mathematikunterricht mit Technologie – ein didaktisches Handbuch mit einer Vielzahl von Aufgaben. \emph{Veritas-Verlag}, Linz, 2014, ISBN 978-3-7101-0431. \newline{}
%  Model for references to proceedings. 

\end{abstract_online}
    